% !TeX root = ../papers/nc-p3.tex
\subsection{有限四辺の全域での添字保持同型と周期}
\label{sec:planar-point-flags}
三次元接点を持たない場合も、局所点環どうしの同型写像をあらかじめ求めずに、
第一・第二窓の添字保持同型を有限個の数値記録で判定できる。

\begin{lemma}\label{lem:planar-point-wronskian}
長さ高々20、埋込み次元2の局所Artin完全交叉 $C$ に対し、
余接平面の直線からなる非空な有限集合が同型不変に定まり、その個数は高々6である。
\end{lemma}
\begin{proof}
$V=\mathfrak m/\mathfrak m^2$ とし、$\Sym V\to\operatorname{gr}_{\mathfrak m} C$ の
核が初めて非零となる次数を $d$、その次数の核を $J_d$ とする。
極小表示 $C=k[[x,y]]/(f,g)$ を取ると、$f,g$ の位数は $d$ 以上で、
次数 $d$ の初期式で $J_d$ が張られる。特に $1\le\dim J_d\le2$ である。
次数 $2d-2$ までの打切り環の次元は $d(2d-1)$ であり、
その中で両関係の倍数の次元は高々 $2\binom d2=d(d-1)$ である。
従って $\dim C\ge d^2$、すなわち $2\le d\le4$ である。

$J_d$ が一次元ならその生成元 $f_d$ を、二次元なら基底 $f_d,g_d$ の
二元Jacobian $W=(f_d)_x(g_d)_y-(f_d)_y(g_d)_x$ を取る。
後者は独立な同次数二元形式の比が標数零で非定数であるため零でない。
例えばアフィン座標で微分して斉次Euler等式を使えばよい。
基底変更と変数変更に対してこの式の射影類は自然に移る。
次数は $d$ または $2d-2$ なので、$\mathbb P(V^*)$ 上の相異なる根の数は1から6である。
各根の零化直線を $V$ 内に取れば主張を得る。
\end{proof}

\begin{theorem}\label{thm:finite-all-window-keys}
四辺が有限な任意の二つの標準AS対象について、添字保持同型の存在と最小周期は、
連続する三つの四次元座標空間からの有限な記録の比較で決定できる。
埋込み次元2以下の部分では、点環の同型型、三空間の独立な単元倍、共通の
局所パラメータ変更を同じ計算に含める。使うのは高々四次の二元形式または
その次数6以下のJacobianの根、因子の有限置換、線型代数および整数核の単項式である。
三次元接点の部分は定理\ref{thm:maximal-tangent-complete}の表示で決定される。
\end{theorem}
\begin{proof}
点環を $D=\prod C_a$ とし、局所因子の長さ・埋込み次元 $e_a$・
$\mathfrak m_a^{n_a}=0$ となる最小の $n_a$ を記録する。
データが等しい因子に標識を付け、$e_a=2$ の各因子で前補題の直線を一つ選ぶ。
適合する余接基底を $x,y$（指定直線 $ky$）とし、これを極大イデアルに持ち上げる。
$e_a=1$ では一変数を取り、$e_a=0$ では変数は不要である。
単項式を評価して得る全射を
\[
 T=\prod_a k[t_1,\ldots,t_{e_a}]/(t_1,\ldots,t_{e_a})^{n_a}
            \longrightarrow D
\]
とし、核を $I$ とする。$U_s\subset D$ を定理\ref{thm:finite-square-chart}の
座標空間とし、その全逆像 $P_s\subset T$ を取る。比較に使うのは
四つの部分空間 $(I,P_0,P_1,P_2)$ である。
この形なら三つの自明化に共通して同じ点環を用いるという条件が失われない。

各因子の単項式を全次数昇順、同じ次数では $y$ の指数昇順に並べる。
標識と直線を保つ任意の点環同型は、$T$ の因子ごとの変数置換に持ち上がる。
実際、選んだ生成元の像の代表多項式を取れば、一次係数は可逆で
$x\mapsto a x+b y$, $y\mapsto c y$ の形を持ち、打切り環の自己同型となる。
逆に $I$ を $I'$ に運ぶ置換は点環同型を与える。任意の局所単元も単元へ持ち上がるので、
三空間の同値は $T$ のこの共通変数変更と $P_s$ ごとの独立な単元倍に等しい。
定数因子のスカラーと、各 $a,c$ をトーラスに分けると、残る群は下単三角である。
そのLie環は次の操作で張られる。各位置別の単元の冪零部分による乗法、
全四空間に共通の次数2以上の単項式を係数とする偏微分、
および二変数因子の $y\partial_x$ である。
いずれも上の行順で次数を上げるか、同次数で $y$ の指数を上げる。
定数項1の単元と指定した線型部の変数変更はその係数で連結にパラメータ化される。

以下は定理\ref{thm:curvilinear-window-orbits}の三角消去と同じであるが、
全逆像の列数は4には限らない。各空間の独立な列基底を取り、
最初の独立な主元行を単位行列に正規化する。非主元行を昇順、列を降順に読み、
既読値の安定化Lie環の平行移動速度の核を順次取る。
速度非零の座標を指数作用で零にすると、単冪軌道の代表は一意である。
この証明は任意の列数で成立し、核 $I$ の次元が零の場合はその行列を空にしてよい。
行 $x^iy^j$ のトーラス重みは $a^ic^j$、単元倍を許す側では
さらにその位置・因子の定数倍を掛ける。主元行で割った重みの整数核を使って
非零値の不変単項式を取る。主元・台・単項式値を全て記録すれば
標識と指定直線を保つ同値と記録の一致が必要十分になる。

従って根と同じ型の因子の全ての標識付けにわたる記録集合が完全不変量となる。
局所完全交叉であることは有限な辺の完全交叉性から従う。
埋込み次元3を含む場合はその存在自体が不変で、前節までの分類を使う。
最小周期は実際の射影対応に沿って得た開始位置 $0,1,2$ の記録を比較して
系\ref{cor:curvilinear-window-period}と同じく決まる。
任意の三空間から作った二窓がAS延長を持つかどうかは、この同型判定とは別の条件である。
\end{proof}

\path{computations/verify_finite_planar_window_orbits.py} では
局所点環 $k[x,y]/(x^2,y^2)$ と16個の単純点の直積を使い、二元Jacobian、
非線型な共通変数変更・三位置の単元倍・核を足した持上げの変更が記録を保つことを検査する。
一位置だけで接係数を変えると標識付き記録が異なることも検査する。
この三空間は軌道計算の試験入力である。
